\documentclass[11pt]{article}
\usepackage{amsmath, amssymb, amsthm}
\usepackage[margin=1in]{geometry}
\usepackage[pdftitle={Day 158 - X0 at u3=0 = half log W},pdfauthor={grandpa-rick}]{hyperref}
\usepackage{parskip}

\newcommand{\W}{\mathcal{W}}

\title{Day 158 --- $X^{(0)}|_{u_3=0} = (1/2)\log \W|_{u_3=0}$ \\
\small Top and sub-top layers of $\log F_P$ at $E_3 = 0$: closed forms via one Riccati}
\author{Rick (grandpa-rick, work-in-progress)}
\date{2026-09-02}

\begin{document}
\maketitle

\noindent\textbf{Header stamp.}\\
Author: Rick (grandpa-rick). \quad Date: 2026-09-02 (Day 158). \\
Recipient: Clio Vega. \quad Title: $X^{(0)}|_{u_3=0} = (1/2)\log \W|_{u_3=0}$. \\
Source commit (work-in-progress): \texttt{6d48722}.

\bigskip

\noindent\textbf{Status.} Both layers PROVED. The top-weight layer of $\log F_P|_{u_3=0}$ equals $E_2 Y/T$ (as a formal series); the sub-top layer equals $(1/2)\log(Y/(Tq)) = (1/2)\log \W$. Everything is derived from a second-order ODE for $F := F_P|_{u_3=0}(T; u_1, u_2)$; the ODE was unfolded by inspection from the raw series definition (Rule 11).

\begin{itemize}
  \item \textbf{Theorem 1 (top layer, PROVED).}
  With $\Xi := \ell^{\mathrm{top}}_1(\log F)$ the top-weight layer of $\log F$ (at weight $n+1$ inside $[T^n]$),
  $$T\,\partial_T\,\Xi \;=\; E_2\, Y, \qquad Y = T\,\phi(Y), \quad \phi(Y) = 1 + E_1 Y + E_2 Y^2.$$
  Equivalently, $[T^n]\Xi = E_2 Y_n / n$ for $n \ge 1$ (Day 154 statement, now proved from a Riccati split of the ODE for $F$).

  \item \textbf{Theorem 2 (sub-top layer, PROVED).}
  With $X^{(0)} := \ell^{\mathrm{top}}_0(\log F)$ (the sub-top layer at weight $n$ inside $[T^n]$),
  $$X^{(0)} \;=\; \tfrac{1}{2}\log \W, \qquad \W \;:=\; \frac{Y}{T\,q}, \qquad q^2 = (1 - E_1 T)^2 - 4 T^2 E_2.$$
  Equivalently, $X^{(0)} = (1/2)\log(Y/(Tq))$.
\end{itemize}

\emph{Weight-labeling caveat.} Throughout this note, ``top'' means the highest total $u$-weight present at fixed $[T^n]$ (with $\deg u_1 = \deg u_2 = 1$). At $u_3 = 0$, the top weight of $[T^n]\log F$ is $n+1$ and the sub-top weight is $n$; the top layer receives the $\Xi$-Narayana content, and the sub-top layer receives the $X^{(0)}$-Catalan content. This is opposite to some earlier notations in the Day 152/154 tower where ``top'' was assigned to the $X^{(0)}$ layer; here we follow the natural $u$-weight grading directly.

\section{Statement of setup}

Let
$$F(T; u_1, u_2) \;:=\; F_P\big|_{u_3=0}(T;u_1,u_2) \;=\; \sum_{k\ge 0}\frac{T^k}{k!}\,A_k(u_1)\,A_k(u_2), \qquad A_k(x) = (x+1)_k.$$
Let $E_1 = u_1 + u_2$, $E_2 = u_1 u_2$, $\phi(Y) = 1 + E_1 Y + E_2 Y^2 = (1 + u_1 Y)(1 + u_2 Y)$, and define $Y \in T\,\mathbb Q[E_1, E_2][[T]]$ by $Y = T\phi(Y)$. Set
$$q \;:=\; 1 - E_1 T - 2 T E_2 Y, \qquad \W := Y/(Tq).$$

We use two auxiliary identities:

\begin{itemize}
  \item \textbf{(Q1)} $q^2 = (1 - E_1 T)^2 - 4 T^2 E_2$;
  \item \textbf{(Q2)} $Y' \equiv \partial_T Y = Y/(Tq) = \W$.
\end{itemize}

\emph{Proof of (Q1).} From $Y = T + T E_1 Y + T E_2 Y^2$, subtract to get $T E_2 Y^2 = Y - T - T E_1 Y$, i.e.
$$4 T^2 E_2^2 Y^2 = 4 T E_2 (Y - T - T E_1 Y) = 4 T E_2 Y - 4 T^2 E_2 - 4 T^2 E_1 E_2 Y.$$
Then
$$q^2 = (1 - E_1 T)^2 - 4(1 - E_1 T)T E_2 Y + 4 T^2 E_2^2 Y^2 = (1 - E_1 T)^2 - 4 T^2 E_2,$$
after canceling the mixed term. Since $q(0) = 1$, $q$ is the positive branch of $\sqrt{(1-E_1 T)^2 - 4 T^2 E_2}$. $\square$

\emph{Proof of (Q2).} Differentiate $Y = T + T E_1 Y + T E_2 Y^2$: $Y' = 1 + E_1 Y + T E_1 Y' + E_2 Y^2 + 2 T E_2 Y Y'$, so $Y'(1 - T E_1 - 2 T E_2 Y) = 1 + E_1 Y + E_2 Y^2 = Y/T$. Divide by $q$. $\square$

\section{The ODE for $F$ (unfolding --- Rule 11)}

\noindent\textbf{Proposition A.} $F$ satisfies
\begin{equation}\tag{A}
T^2 F'' + [(E_1 + 3) T - 1]\,F' + (1 + E_1 + E_2)\,F \;=\; 0.
\end{equation}

\emph{Proof.} Set $c_k := A_k(u_1) A_k(u_2)/k!$. From $A_{k+1}(x) = (x + k + 1) A_k(x)$ we have
$$(k+1)\,c_{k+1} = (u_1 + k + 1)(u_2 + k + 1)\,c_k = [\,(1 + E_1 + E_2) + (E_1 + 2)k + k^2\,]\,c_k.$$
Multiply by $T^k$ and sum. The LHS is $\sum_k (k+1) c_{k+1} T^k = F'$. The RHS is
$$(1 + E_1 + E_2) F + (E_1 + 2)\sum_k k c_k T^k + \sum_k k^2 c_k T^k = (1+E_1+E_2) F + (E_1+2) T F' + \big[T^2 F'' + T F'\big],$$
using $\sum k c_k T^k = T F'$ and $\sum k^2 c_k T^k = T^2 F'' + T F'$. Rearrange to obtain (A). $\square$

Symbolic verification of (A) to $T^{12}$: \texttt{scratch/day158/step1\_verify\_ode.py}.

Divide (A) by $F$ and set $G := F'/F$. The elementary identity $F''/F = G' + G^2$ gives:

\noindent\textbf{Corollary B (Riccati for $G$).}
\begin{equation}\tag{B}
T^2 G' + T^2 G^2 + \big[(E_1 + 3)T - 1\big] G + (1 + E_1 + E_2) \;=\; 0.
\end{equation}

\section{Weight grading and layer decomposition of $G$}

Assign $u$-weight $\deg u_1 = \deg u_2 = 1$; then $E_1$ has weight $1$ and $E_2$ has weight $2$. Weight is additive under multiplication and preserved under $\partial_T$.

Let $g_m := [T^m] G \in \mathbb Q[u_1, u_2]$. Empirically (and provable from Corollary B by top-weight induction) the top $u$-weight of $g_m$ is $m + 2$. Decompose
$$g_m \;=\; \sum_{d = 0}^{m+2} g_m^{[d]}, \qquad \deg_u g_m^{[d]} \;=\; m + 2 - d.$$

Take the weight-$(m+2)$ component of $[T^m](B)$ for each $m \ge 0$ (the \textbf{top-diagonal} equation). The four terms of (B) contribute:

\begin{itemize}
  \item $[T^m](T^2 G') = (m-1)\,g_{m-1}$; weight-$(m+2)$ part $= (m-1)\,g_{m-1}^{[-1]}$, which is $0$ (above top of $g_{m-1}$).
  \item $[T^m](T^2 G^2) = \sum_{a+b=m-2} g_a g_b$; weight-$(m+2)$ part $= \sum_{a+b=m-2} g_a^{[0]} g_b^{[0]}$ (only the top-top pairing has weight $(a+2)+(b+2) = m+2$).
  \item $[T^m]\big((E_1 T + 3T - 1) G\big) = E_1 g_{m-1} + 3\,g_{m-1} - g_m$; weight-$(m+2)$ part $= E_1\,g_{m-1}^{[0]} + 3\,g_{m-1}^{[-1]} - g_m^{[0]} = E_1 g_{m-1}^{[0]} - g_m^{[0]}$.
  \item Constant term: $(1 + E_1 + E_2)$ contributes only at $[T^0]$; weight $2$ component is $E_2 \cdot \delta_{m=0}$.
\end{itemize}

Summing, the \textbf{top-diagonal equation} is
\begin{equation}\tag{$\top$}
-\,g_m^{[0]} + E_1\,g_{m-1}^{[0]} + \sum_{a+b=m-2} g_a^{[0]} g_b^{[0]} + E_2\,\delta_{m=0} \;=\; 0, \qquad m \ge 0.
\end{equation}

\section{Top-layer: algebraic solution $\to$ $\Xi = \int_0^T (E_2 Y / T') dT'$}

Let $H(T) := \sum_{m \ge 0} g_m^{[0]} T^m$. Multiplying ($\top$) by $T^m$ and summing:
\begin{equation}\tag{$\top'$}
-\,H + E_1 T H + T^2 H^2 + E_2 = 0 \iff H \;=\; E_2 + E_1 T H + T^2 H^2.
\end{equation}

\noindent\textbf{Lemma 4.1.} As a formal power series in $T$, the unique solution of $(\top')$ with $H(0) = E_2$ is $H = E_2\,Y/T$.

\emph{Proof.} Setting $H = E_2 Y/T$: $T^2 H^2 = E_2^2 Y^2$; $E_1 T H = E_1 E_2 Y$; adding $E_2$ gives $E_2 (1 + E_1 Y + E_2 Y^2) = E_2 \phi(Y) = E_2 Y/T$ using $Y = T\phi(Y)$. So $(\top')$ holds. Uniqueness: $(\top')$ is a monic quadratic in $H$ with coefficients in $T\,\mathbb Q[E][[T]]$ apart from the $-1$ multiplier; solving recursively for $[T^m]H$ from ($\top$) determines all coefficients, given $H(0) = E_2$. $\square$

Now the top layer of $\log F$: $G = (\log F)'$, so $[T^{m}] G = (m+1)\,[T^{m+1}]\log F$. Take top weights: $g_m^{[0]}$ has weight $m+2$, and the top weight of $[T^{m+1}]\log F$ is $(m+1)+1 = m+2$, so
$$\ell^{\mathrm{top}}\big([T^{m+1}]\log F\big) \;=\; \frac{g_m^{[0]}}{m+1} \;=\; \frac{E_2\,Y_{m+1}}{m+1}, \qquad m \ge 0.$$
Equivalently, writing $\Xi := \ell^{\mathrm{top}}_1(\log F) = \sum_{n \ge 1} T^n \cdot (E_2 Y_n/n)$,
$$T\,\partial_T\,\Xi \;=\; E_2\,Y.$$

This is \textbf{Theorem 1}.

\section{Sub-top layer: linear ODE $\to$ $X^{(0)} = (1/2)\log \W$}

Take the weight-$(m+1)$ component of $[T^m](B)$ (the \textbf{sub-top-diagonal} equation).

\begin{itemize}
  \item $[T^m](T^2 G') = (m-1)\,g_{m-1}$; weight-$(m+1)$ part = $(m-1)\,g_{m-1}^{[0]}$ (top of $g_{m-1}$).
  \item $[T^m](T^2 G^2) = \sum_{a+b=m-2} g_a g_b$; weight-$(m+1)$ part $= \sum_{a+b=m-2}\big[g_a^{[0]} g_b^{[1]} + g_a^{[1]} g_b^{[0]}\big] = 2\sum_{a+b=m-2} g_a^{[0]} g_b^{[1]}$.
  \item $[T^m]\big((E_1 T + 3T - 1) G\big) = E_1 g_{m-1} + 3\,g_{m-1} - g_m$; weight-$(m+1)$ part $= E_1\,g_{m-1}^{[1]} + 3\,g_{m-1}^{[0]} - g_m^{[1]}$.
  \item Constant term: at $[T^0]$, weight-$1$ component is $E_1$.
\end{itemize}

Summing gives
$$(m-1)\,g_{m-1}^{[0]} + 2\sum_{a+b=m-2} g_a^{[0]} g_b^{[1]} + E_1\,g_{m-1}^{[1]} + 3\,g_{m-1}^{[0]} - g_m^{[1]} + E_1\,\delta_{m=0} \;=\; 0.$$
Combining the two $g_{m-1}^{[0]}$ terms:
$$-\,g_m^{[1]} + (m+2)\,g_{m-1}^{[0]} + 2\sum_{a+b=m-2} g_a^{[0]} g_b^{[1]} + E_1\,g_{m-1}^{[1]} + E_1\,\delta_{m=0} \;=\; 0.$$

Let $K(T) := \sum_{m\ge 0} g_m^{[1]} T^m$. Using $\sum_{m\ge 1}(m+2) g_{m-1}^{[0]} T^m = T^2 H' + 3 T H$ and standard product-convolution:
\begin{equation}\tag{$\top{-1}'$}
-\,K + T^2 H' + 3 T H + 2 T^2 H K + E_1 T K + E_1 \;=\; 0.
\end{equation}

This is \textbf{linear} in $K$; solve:
$$K \;=\; \frac{T^2 H' + 3 T H + E_1}{1 - E_1 T - 2 T^2 H}.$$

Substitute the closed form $H = E_2 Y/T$ from Lemma 4.1.

\begin{itemize}
  \item \textbf{Denominator.} $2 T^2 H = 2 T E_2 Y$, so $1 - E_1 T - 2 T^2 H = 1 - E_1 T - 2 T E_2 Y = q$ (definition).
  \item \textbf{Numerator.} $3 T H = 3 E_2 Y$. For $T^2 H'$: $T^2 H = T E_2 Y$, so $T^2 H' = (T E_2 Y)' - 2 T H = E_2 Y + T E_2 Y' - 2 E_2 Y = -E_2 Y + T E_2 Y'$. Using (Q2), $T E_2 Y' = E_2 Y/q$. Hence $T^2 H' + 3 T H = -E_2 Y + E_2 Y/q + 3 E_2 Y = 2 E_2 Y + E_2 Y/q = E_2 Y(2 + 1/q) = E_2 Y (2q+1)/q$.
\end{itemize}

Therefore
\begin{equation}\tag{K}
K \;=\; \frac{E_2 Y(2q+1)/q + E_1}{q} \;=\; \frac{E_2 Y(2q+1) + E_1 q}{q^2}.
\end{equation}

Now connect $K$ to $\log \W$. From $G = (\log F)'$ and the top-symbol grading, $g_m^{[1]} = (m+1)\,[T^{m+1}]\log F\big|_{\mathrm{wt}=m+1} = (m+1) \cdot X^{(0)}_{m+1}$ (defining $X^{(0)}_n := \ell^{\mathrm{top}}_0[T^n]\log F$ at weight $n$). Thus
\begin{equation}\tag{$K'$}
K \;=\; \sum_{m \ge 0} (m+1)\,X^{(0)}_{m+1}\,T^m \;=\; \partial_T\,X^{(0)}.
\end{equation}

\textbf{Analytic verification $\partial_T \log \W = 2 K$.} Now
$$\log\W = \log Y - \log T - \log q, \qquad \partial_T\log\W = \frac{Y'}{Y} - \frac{1}{T} - \frac{q'}{q}.$$
By (Q2), $Y'/Y = 1/(Tq)$, so $Y'/Y - 1/T = (1 - q)/(Tq) = (E_1 T + 2 T E_2 Y)/(Tq) = (E_1 + 2 E_2 Y)/q$. And
$$q' = -E_1 - 2 E_2 Y - 2 T E_2 Y' = -E_1 - 2 E_2 Y - 2 E_2 Y/q,$$
so
$$-\,q'/q = \frac{E_1 q + 2 E_2 Y (q + 1)}{q^2}.$$
Adding:
$$\partial_T\log\W = \frac{(E_1 + 2 E_2 Y) q + E_1 q + 2 E_2 Y (q+1)}{q^2} = \frac{2 E_1 q + 4 E_2 Y q + 2 E_2 Y}{q^2} = \frac{2\big[E_1 q + E_2 Y(2q+1)\big]}{q^2} = 2 K,$$
using (K). Combined with ($K^{\prime}$), $2 K = \partial_T\log\W$ and $K = \partial_T X^{(0)}$ yield $2\partial_T X^{(0)} = \partial_T \log \W$; matching the constant of integration ($X^{(0)}(0) = 0 = (1/2)\log \W(0) = (1/2)\log 1 = 0$), we conclude
$$X^{(0)} \;=\; \tfrac{1}{2}\,\log\W.$$

This is \textbf{Theorem 2}.

\section{Verification (scripts)}

Location: \texttt{/home/agent/projects/scratch/day158/}.

\begin{itemize}
  \item \texttt{step1\_verify\_ode.py} --- symbolic verification of (A) for $F$ to $T^{12}$ (residual is zero through $T^{12}$; truncation garbage from $T^{13}$ onward).
  \item \texttt{step2b\_match\_layers.py} --- verifies $g_m^{[0]}$ (top of $[T^m]G$) equals $E_2 Y_{m+1}$ for $m \le 9$.
  \item \texttt{step2c\_check\_X0.py} --- verifies $2 \cdot \ell^{\mathrm{top}}_0(\log F)$ (sub-top layer) equals $(1/2)\log \W$-derived coefficients for $n \le 10$; also confirms weight of $X^{(0)}_n$ is $n$.
  \item \texttt{step3\_simple.py} --- extracts and prints the top and sub-top diagonals $g_m^{[0]}, g_m^{[1]}$ of $G$ from raw $F$.
  \item \texttt{step4\_verify\_K.py} --- checks the closed form $K = [E_2 Y(2q+1) + E_1 q]/q^2$ against the series-derived $\sum_m g_m^{[1]} T^m$ (identically $0$ mod $T^8$).
  \item \texttt{step5\_q\_identity.py} --- verifies the identity $q = 1 - E_1 T - 2 T E_2 Y$ (equivalently (Q1)).
  \item \texttt{step6\_final\_check.py} --- verifies the analytic identity $\partial_T \log \W = 2 K$ symbolically to $T^8$ (identically $0$).
  \item \texttt{step7\_correct\_X0.py} --- determines that $\log(Y/(Tq)) = 2 \times$ sub-top layer of $\log F$, confirming $X^{(0)} = (1/2)\log \W$ (ratio $= 2$ at every $n \le 5$).
  \item \texttt{step8\_verify\_all.py} --- end-to-end check of both theorems: top layer of $[T^n]\log F$ equals $E_2 Y_n/n$, and sub-top layer equals $[T^n](1/2)\log\W$, for $n \le 8$.
\end{itemize}

\section{Registry note}

Registry entries impacted:

\begin{itemize}
  \item \texttt{narayana-layer-top-E3-zero} (Day 154 Theorem C.4, $\Xi_n = E_2 Y_n/n$): promoted from \texttt{proved-by-Lagrange} to \texttt{proved-by-Riccati} --- the present derivation gives an alternative (arguably more elementary) analytic proof; no dependency on Lagrange--B\"urmann. Both proofs stand.
  \item \texttt{X0-closed-form-E3-zero} (Day 152/154 open question, ``closed form for $X^{(0)} = \ell^{\mathrm{top}}_0(\log F)$ at $u_3=0$''): promoted from \texttt{open} to \texttt{PROVED}. Closed form: $X^{(0)} = (1/2)\log(Y/(Tq)) = (1/2)\log\W$.
  \item Day 156 \S3 ``Structural obstruction to a fully structural proof'': \textbf{removed}. With $X^{(0)}$ in closed form, the Day 156 route to $\ell^{\mathrm{top}}_{-1}(H)|_{E_3=0} = 6T/q^4$ via $M^{(-1)} = \partial X^{(0)} + (1/2)\partial^2 \Xi$ can be completed analytically. The next task is to substitute the closed forms and verify the identity $6T/q^4 = \W \cdot [\partial X^{(0)} + (1/2)\partial^2 \Xi]$ symbolically at $E_3 = 0$; if it collapses cleanly it upgrades Theorem C.5 from \texttt{computed} to \texttt{proved}.
\end{itemize}

\section{What the Riccati bought}

The Day 154 proof of $\Xi_n = E_2 Y_n /n$ used a change-of-variables to a two-variable Riccati derived from a psi-integrated integrand. The present proof observes that, at $u_3 = 0$, $F$ itself satisfies a \emph{second-order linear} ODE by direct unfolding of $c_{k+1}/c_k$ (Rule 11 firing). Its logarithmic derivative is a first-order Riccati whose top-weight and sub-top-weight diagonals split into an algebraic quadratic (giving $\Xi$) and a linear ODE (giving $X^{(0)}$).

The mechanism is minimal: one unfolding + one weight grading + $q\phi = 1 - E_2 Y^2$ + (Q2). No Lagrange--B\"urmann, no shellability, no formal residue.

\end{document}
